\begin{Exercise}[label={ex:lettura_scrittura}]
    Dato un file di testo "people.txt" che contiene i dati di una serie di persone (non più di 20),
    una persona per riga.
    Si scriva un programma che, una volta letti da file i dati di queste persone ne estragga l'insieme di persone
    compatibili con una nuova persona data salvando il risultato sul file binario "partners.dat".
    Due persone sono compatibili se sono di sesso diverso e la differenza di età, riferita solo all'anno, non supera i 5 anni.

    Ogni riga del file di testo "people.txt" ha il seguente formato:
    \begin{itemize}
        \item cognome (stringa senza spazi, al massimo 19 caratteri).
        \item un separatore ';'.
        \item nome (stringa senza spazi, al massimo 19 caratteri).
        \item un separatore ';'.
        \item data di nascita (formato gg/mm/aaaa).
        \item un solo spazio.
        \item un carattere che indica il sesso ('M' o 'F').
    \end{itemize}

    Si definiscano le opportune strutture dati e si realizzino le seguenti funzioni:
    \begin{itemize}
        \item una funzione \texttt{lettura()} che dato il nome del file, legge i dati delle persone
        e li memorizza in un array di nome \texttt{elenco}.
        \item una funzione \texttt{compatibili()} che date due persone restituisce 1 solo se le due persone sono
        compatibili.
        \item un \texttt{main()} che invochi la lettura per acquisire le persone, poi chieda all'utente i dati
        di una nuova persona per scrivere sul file binario l'insieme delle persone compatibili con essa.
    \end{itemize}
\end{Exercise}

\begin{Answer}[ref={ex:lettura_scrittura}]
  	\lstinputlisting[style=cstylesmall]{src/lettura_scrittura.c}
\end{Answer}